Peripersonal space (PPS) is often inferred from tactile reaction time during an approaching sound. We tested whether scheduled respiratory phase changed this SOA-dependent response differently in breath-coupled virtual reality (Viscereality) and a control condition. The cleared pool contained 40 participants and 8{,}495 retained audio--tactile hit trials. The preregistered Condition $\times$ Phase $\times$ SOA mixed-model interaction was not detected, $\chi^{2}(1)=0.006$, $p=.939$; categorical-SOA and partial-session-exclusion sensitivities agreed. Audio--tactile responses were 17.9~ms faster than tactile-only Baseline responses on average, but source recordings from 15 participants showed that their tactile cues occurred 500~ms later than Baseline cues at the same nominal SOA label, limiting Baseline subtraction. None of 160 participant-by-condition-by-phase trial-level sigmoid fits met the registered $R^{2}\geq .70$ threshold. Group curves were descriptive and did not yield defensible participant-level PPS midpoints or sharpness estimates. Baseline reaction time varied with scheduled phase. These data do not demonstrate the predicted Viscereality-specific respiratory modulation. Differences in source reaction-time definitions and the lack of a uniform recording-based respiratory-phase audit further limit physiological interpretation.
